

\section{Experiments and Analysis}
\label{sec:experiments}

\subsection{Experimental Setup}
\label{sec:experimental-setup}



\paragraph{Evaluated models.}
We evaluate DME at two model scales, denoted as DME-2B and DME-9B. Both models follow the same two-stage training framework described in Section~\ref{sec:multi-stage-learning-framework}. Stage 1 performs large-scale contrastive pre-training, while Stage 2 further applies Evidence-Grounded Typed Latent Reasoning and Cross-Conditional Reconstruction. Unless otherwise specified, the final DME models use the complete Stage-2 training recipe.

\paragraph{Benchmarks.}
For public evaluation, we report results on MMEB-v2, a comprehensive benchmark for universal multimodal embedding that covers diverse retrieval tasks across text, image, video, visual document, and mixed-modality inputs~\cite{MMEB,VLM2VecV2,Qwen3VLEmbedding}. We use MMEB-v2 as the main benchmark to compare DME against existing multimodal embedding models under comparable model-scale regimes. For internal evaluation, we additionally use Douyin's in-house offline evaluation set, which reflects real-world production retrieval traffic across cross-modal directions, to assess whether the DME techniques transfer to a large-scale industrial setting.


\paragraph{Baselines.}
We compare DME with representative MLLM-based multimodal embedding models and recent reasoning-enhanced embedding methods. The MLLM-based embedding baselines include GME~\cite{GME}, VLM2Vec/VLM2Vec-V2~\cite{MMEB,VLM2VecV2}, and Qwen3-VL-Embedding~\cite{Qwen3VLEmbedding}. We also compare with recent reasoning-oriented retrieval representation methods, including Think-Then-Embed (TTE)~\cite{ThinkThenEmbed}, TTE-v2~\cite{ThinkThenEmbedV2}, and Embed-RL~\cite{EmbedRL}. When reporting scale-wise comparisons, we group baselines by model size to ensure that DME-2B and DME-9B are compared with models of comparable capacity.

\paragraph{Metrics.}
Following the standard MMEB-v2 evaluation protocol, we report the official aggregate score together with task-level or modality-level breakdowns when available. For retrieval tasks, we use ranking-based metrics such as Recall@$K$, nDCG@$K$, or task-specific official metrics depending on the benchmark definition. For internal Douyin evaluation, we report the relative improvement of DME over the in-production baseline on Douyin's in-house offline evaluation set.

\begin{table*}[t]
\centering
\small
\setlength{\tabcolsep}{3.1pt}
\renewcommand{\arraystretch}{1.08}
\resizebox{\linewidth}{!}{
\begin{tabular}{l c ccccc ccccc ccccc c}
\toprule
\multirow{2}{*}{\textbf{Model}} &
\multirow{2}{*}{\textbf{Size}} &
\multicolumn{5}{c}{\textbf{Image}} &
\multicolumn{5}{c}{\textbf{Video}} &
\multicolumn{5}{c}{\textbf{VisDoc}} &
\multirow{2}{*}{\textbf{All}} \\
\cmidrule(lr){3-7}
\cmidrule(lr){8-12}
\cmidrule(lr){13-17}
& & CLS & QA & RET & GD & Avg.
& CLS & QA & RET & MRET & Avg.
& VDRv1 & VDRv2 & VR & OOD & Avg.
& \\
\midrule
\# datasets
& -- & 10 & 10 & 12 & 4 & 36
& 5 & 5 & 5 & 3 & 18
& 10 & 4 & 6 & 4 & 24
& 78 \\
\midrule
VLM2Vec~\cite{MMEB} & 8B
& 62.7 & 56.9 & 69.4 & 82.2 & 65.5
& 39.1 & 30.0 & 29.0 & 38.9 & 33.7
& 56.9 & 9.4 & 59.1 & 54.0 & 49.1
& 53.1 \\
VLM2Vec-V2~\cite{VLM2VecV2} & 2B
& 62.9 & 56.3 & 69.5 & 77.3 & 64.9
& 39.3 & 34.3 & 28.8 & 36.8 & 34.6
& 75.5 & 44.9 & 79.4 & 62.2 & 69.2
& 59.2 \\
GME~\cite{GME} & 8B
& 57.7 & 34.7 & 71.2 & 59.3 & 56.0
& 37.4 & 50.4 & 28.4 & 37.0 & 38.4
& 89.4 & 55.6 & 85.0 & 68.3 & 79.3
& 59.1 \\
Ops-MM-Embedding-v1\modelhome{https://huggingface.co/OpenSearch-AI/Ops-MM-embedding-v1-7B} & 8B
& 69.7 & 69.6 & 73.1 & 87.2 & 72.7
& 59.7 & 62.2 & 45.7 & 43.2 & 53.8
& 80.1 & 59.6 & 79.3 & 67.8 & 74.4
& 68.9 \\
RzenEmbed~\cite{RzenEmbed}\modelhome{https://github.com/360CVGroup/RzenEmbed} & 8B
& 70.6 & 71.7 & 78.5 & 92.1 & 75.9
& 58.8 & 63.5 & 51.0 & 45.5 & 55.7
& 89.7 & 60.7 & 88.7 & 69.9 & 81.3
& 72.9 \\

Embed-RL~\cite{EmbedRL} & 4B
& 63.7 & 70.5 & 71.3 & 91.4 & 70.1
& 57.6 & 58.4 & 45.1 & 49.5 & 53.0
& 80.2 & 53.4 & 84.9 & 67.1 & 74.7
& 68.1 \\

WeMM-Embedding & 2B
& 72.1 & 72.6 & 76.6 & 93.3 & 76.1
& 61.5 & 64.0 & 54.2 & 52.6 & 58.7
& 87.2 & 53.4 & 88.8 & 34.5 & 73.2
& 71.2 \\

Qwen3-VL-Embedding~\cite{Qwen3VLEmbedding}\modelhome{https://huggingface.co/Qwen/Qwen3-VL-Embedding-2B} & 2B
& 70.3 & 74.3 & 74.8 & 88.5 & 75.0
& 71.9 & 64.9 & 53.9 & 53.3 & 61.9
& 84.4 & 65.3 & 86.4 & 69.4 & 79.2
& 73.2 \\
IFM-TTE~\cite{ThinkThenEmbed}\modelhome{https://interestfm-tte.github.io/} & 8B
& 76.7 & 78.5 & 74.6 & 89.3 & 77.9
& 60.5 & 67.9 & 51.7 & 54.9 & 59.2
& 85.2 & 71.5 & 92.7 & 53.3 & 79.5
& 74.1 \\

TTE-v2$^{*}$~\cite{ThinkThenEmbedV2} & 7B
& 78.1 & 79.0 & 76.3 & 91.6 & 79.2
& 58.3 & 66.9 & 52.3 & 68.0 & 60.7
& 85.4 & 63.3 & 94.3 & 69.0 & 82.3
& 75.7 \\

WeMM-Embedding & 8B
& 73.6 & 76.1 & 78.6 & 92.9 & 78.1
& 66.5 & 71.7 & 56.4 & 55.2 & 63.2
& 89.5 & 59.3 & 90.4 & 35.1 & 75.6
& 73.9 \\

Qwen3-VL-Embedding~\cite{Qwen3VLEmbedding}\modelhome{https://huggingface.co/Qwen/Qwen3-VL-Embedding-8B} & 8B
& 74.2 & 81.1 & 80.2 & 92.3 & 80.1
& 78.4 & 71.0 & 58.7 & 56.1 & 67.1
& 87.2 & 69.9 & 88.7 & 73.3 & 82.4
& 77.8 \\
\midrule
\textbf{DME} & \textbf{2B}
& \textbf{68.6} & \textbf{76.2} & \textbf{75.6} & \textbf{94.3} & \textbf{75.9}
& \textbf{84.5} & \textbf{61.9} & \textbf{55.5} & \textbf{57.4} & \textbf{65.6}
& \textbf{87.0} & \textbf{56.2} & \textbf{90.3} & \textbf{70.0} & \textbf{79.9}
& \textbf{74.8} \\
\textbf{DME} & \textbf{9B}
& \textbf{74.5} & \textbf{80.9} & \textbf{78.2} & \textbf{94.6} & \textbf{79.8}
& \textbf{87.7} & \textbf{71.0} & \textbf{61.0} & \textbf{58.5} & \textbf{70.8}
& \textbf{87.6} & \textbf{57.8} & \textbf{94.5} & \textbf{73.5} & \textbf{82.0}
& \textbf{78.4} \\
\bottomrule
\end{tabular}
}
\caption{
\textbf{MMEB-v2 evaluation with task-family breakdown.} CLS, QA, RET, GD, MRET, VDR, VR, and OOD denote classification, question answering, retrieval, grounding, moment retrieval, ViDoRe, VisRAG, and out-of-distribution evaluation, respectively. Unless otherwise noted, all reported scores are taken from the official MMEB Leaderboard. The $\dagger$ marker indicates an available checkpoint, model card, or project homepage. $^{*}$TTE-v2 is reported from its paper under a 76-task MMEB-v2 setting that excludes two VisDoc OOD datasets; other rows follow the 78-task setting when available.
}
\label{tab:mmeb-v2-results}
\end{table*}


\subsection{Results on MMEB-v2}
\label{sec:mmeb-v2-results}

Table~\ref{tab:mmeb-v2-results} reports the main results on MMEB-v2. We compare DME with recent MLLM-based multimodal embedding models and reasoning-enhanced retrieval models under comparable scale regimes. DME-2B obtains an overall score of \textbf{74.8}, and DME-9B obtains an overall score of \textbf{78.4}. From the results, we make three observations.

\paragraph{(1) DME achieves strong scale-wise performance.}
DME-2B should be primarily compared with other 2B-level models such as VLM2Vec-V2 and Qwen3-VL-Embedding-2B, while DME-9B should be compared with larger models such as RzenEmbed-8B, IFM-TTE-8B, TTE-v2-7B, and Qwen3-VL-Embedding-8B. Under this comparison, DME shows strong performance at both model scales, indicating that the proposed two-stage training framework is effective for both compact and larger MLLM-based embedders.

\paragraph{(2) The gains are not concentrated in a single modality group.}
DME performs consistently across Image, Video, and VisDoc groups. For the 2B and 9B variants, DME obtains \textbf{75.9}/\textbf{79.8} on Image, \textbf{65.6}/\textbf{70.8} on Video, and \textbf{79.9}/\textbf{82.0} on VisDoc, respectively. This broad improvement is important because MMEB-v2 contains heterogeneous task families, including classification, question answering, retrieval, grounding, moment retrieval, visual-document retrieval, VisRAG, and out-of-distribution evaluation.

\paragraph{(3) DME outperforms reasoning-enhanced baselines while using lightweight latent-token inference.}
Recent models such as IFM-TTE, TTE-v2, and Embed-RL introduce reasoning signals to improve multimodal retrieval representations. DME follows a different design: evidence grounding and typed latent reasoning are integrated into the encoder through a small number of retrieval-specific latent tokens, instead of generating explicit chain-of-thought text or applying cross-encoder reranking at evaluation time. The results in Table~\ref{tab:mmeb-v2-results} suggest that this latent-token design provides strong retrieval quality while keeping inference close to the standard dense retrieval pipeline.

% \subsection{Results on Douyin Industrial Benchmark}
\label{sec:douyin-industrial-benchmark}

\begin{table}[t]
\centering
\begin{tabular}{lccccc}
\toprule
 & Text2Video & Text2Image & Image2Image & Image2Video & All \\
\midrule
$\Delta$ & +3.10\% & +3.03\% & +2.70\% & +2.83\% & +2.92\% \\
\bottomrule
\end{tabular}
\caption{\textbf{Relative performance gains on Douyin industrial benchmark.} We initialize the in-house model from DME and continue
training with several DME techniques transferred to the in-house setting.
Improvements ($\Delta$) are relative gains over the previous production model.}
\label{tab:inhouse-transfer}
\end{table}

\subsection{Industrial Deployment}
\label{sec:industrial-deployment}

Beyond public benchmarks, we further validate DME in Douyin's production retrieval system through both offline and online evaluations.

\textbf{Offline results.}
We initialize the in-house model from DME and continue training with several DME techniques transferred to this setting. As shown in Table~\ref{tab:inhouse-transfer}, this yields a 2.92\% relative improvement in overall retrieval quality over the previous production model. The gains are consistent across all four directions, ranging from +3.10\% on Text2Video to +2.70\% on Image2Image, which indicates that the DME techniques generalize to a large-scale industrial retrieval setting rather than being effective only under public academic benchmarks.

\textbf{Online results.}
We further deploy the DME-based model in Douyin's online search system, powering scenarios such as generative search and serving as a retrieval feature for ranking. Online A/B testing on Douyin search verifies a 0.1\% Lifetime (LT) gain in core online business metrics.


\subsection{Training Recipe Diagnostics}
\label{sec:training-recipe-diagnostics}

To inspect how each part of the DME learning recipe affects retrieval quality, we analyze the cumulative DME training recipe on MMEB-v2. All variants are trained on the same Stage-2 data and differ only in which components of the recipe are enabled, so the comparison isolates the marginal effect of each component rather than removing modules from the final model. Starting from a model that is trained directly on this data, we progressively enable large-scale Stage-1 pre-training, Stage 2-A, and Stage 2-B, and report how the model changes across all task groups.

We report four configurations. Baseline is initialized from Qwen3.5 and trained directly on the Stage-2 data with the contrastive objective only, without large-scale Stage-1 pre-training and without either Stage-2 mechanism. +Stage 1 adds large-scale contrastive pre-training prior to this training. +Stage 2-A further adds Evidence-Grounded Typed Latent Reasoning, including anchor-based evidence grounding, typed query-side latent supervision,
and evidence-enhanced readout. +Stage 2-B adds Cross-Conditional Reconstruction with NTP/MTP, giving the final DME model.

\begin{table*}[t]
\centering
\setlength{\tabcolsep}{12pt}
% \renewcommand{\arraystretch}{1.08}
\label{tab:ablation}
\begin{tabular}{ccccccc}
\toprule
\multicolumn{3}{c}{Configuration} & \multicolumn{4}{c}{Score} \\
\cmidrule(lr){1-3} \cmidrule(lr){4-7}
Stage 1 & Stage 2-A & Stage 2-B & Image & Video & VisDoc & All \\
\midrule
           &            &            & 74.6 & 55.3 & 77.1 & 70.9 \\
\checkmark &            &            & 74.8 & 59.3 & 79.0 & 72.5 \\
\checkmark & \checkmark &            & 75.2 & 63.7 & 79.2 & 73.8 \\
\checkmark & \checkmark & \checkmark & \textbf{75.9} & \textbf{65.6} & \textbf{79.9} & \textbf{74.8} \\
\bottomrule
\end{tabular}
\caption{
\textbf{Cumulative analysis of the DME training recipe on MMEB-v2.}
All configurations use the Stage-2 contrastive objective; checkmarks indicate
additional training components enabled on top of this common baseline.
Stage 2-A denotes Evidence-Grounded Typed Latent Reasoning, and Stage 2-B denotes
Cross-Conditional Reconstruction. Scores are reported on the Image, Video, and
VisDoc groups together with the overall average.
}
\label{tab:ablation}
\end{table*}

Table~\ref{tab:ablation} summarizes the cumulative effect of the DME recipe, which raises the overall score from 70.9 to 74.8, a gain of 3.9 points. Adding Stage 1 improves the overall score from 70.9 to 72.5, confirming that large-scale heterogeneous contrastive pre-training establishes a stronger unified embedding space even when the same Stage-2 data is used; the gain is concentrated on video (55.3 to 59.3) and visual-document retrieval (77.1 to 79.0), while image retrieval is largely unchanged (74.6 to 74.8). Adding Stage 2-A contributes 1.3 points overall (72.5 to 73.8) and yields the largest single improvement on video (59.3 to 63.7), consistent with its design of grounding retrieval in localized textual, visual, and temporal evidence. Adding Stage 2-B brings the overall score to 74.8, with balanced gains across all three groups (image 75.2 to 75.9, video 63.7 to 65.6, visual-document 79.2 to 79.9), indicating that preserving fine-grained counterpart-side semantics through cross-conditional reconstruction benefits retrieval broadly rather than any single modality.

\subsection{Measuring Semantic Sufficiency via Representation Completeness}
\label{sec:ntp-topk}
To quantify how much content a single embedding actually preserves, we recast
the NTP objective of Eq.~\eqref{eq:s2b-ntp-q2d} as a bounded, interpretable
metric. Instead of the unbounded cross-entropy, we ask whether the trained DME
model can recover each target token from the pooled embedding alone. Consider
the Q$\rightarrow$D direction. We prepend the query prefix
$\tilde{\mathbf{z}}_q$ as the \emph{only} conditioning input, and run a single
teacher-forcing pass through the same backbone over the document-side 
text $x_d=(x_d^1,\dots,x_d^{t})$. At each position, this yields the         
next-token distribution
$P_\theta\!\left(x_d^{t}\mid \tilde{\mathbf{z}}_q, x_d^{<t}\right)$
of Eq.~\eqref{eq:s2b-softmax}. A position counts as a hit if the ground-truth
token lies within the model's Top-$K$ predictions, and we report the
micro-averaged accuracy over the first $L_i$ evaluated positions of each sample:
\begin{equation}
\mathrm{acc}@K \;=\;
\frac{\displaystyle\sum_{i}\sum_{t=1}^{L_i}
      \mathbf{1}\!\left[\,x^{t,(i)}\in
      \mathrm{Top}_K\!\big(P_\theta(\cdot\mid \tilde{\mathbf{z}}^{(i)}, x^{<t,(i)})\big)\right]}
     {\displaystyle\sum_{i} L_i}\;\in\;[0,1], 
\label{eq:acck}
\end{equation}
where $T_i$ is the target length of sample $i$, $L_{\max}$ is a fixed truncation length so that only the first $L_i=\min(L_{\max}, T_i)$ positions of each sample are evaluated (we set $L_{\max}=10$), and $\mathrm{Top}_K(\cdot)$ is the set of $K$ tokens with the largest probability. Unlike the NTP cross-entropy, $\mathrm{acc}@K$ lies in $[0,1]$ and reads directly as the fraction of the first $L_i$ target tokens the embedding can recover within its Top-$K$ guesses, making it comparable across datasets, directions, and
modalities.

We probe four prefix$\rightarrow$target directions per retrieval pair, reusing
the training-time branch construction. The self directions, q2q and d2d,
condition on $\tilde{\mathbf{z}}_q$ (resp.\ $\tilde{\mathbf{z}}_d$) to recover
its own text $x_q$ (resp.\ $x_d$), and thus measure \emph{reconstruction
completeness}: how faithfully an embedding re-encodes its own content. The
cross directions, q2d and d2q, condition on one side to recover the other
(mirroring $\mathcal{L}_{\text{NTP}}^{q\rightarrow d}$ and
$\mathcal{L}_{\text{NTP}}^{d\rightarrow q}$), and thus measure
\emph{counterpart completeness}: how much of the paired document (or query) a
query (or document) embedding retains. A high $\mathrm{acc}@K$ indicates that
the pooled vector exposes, through the DME model, the token-level content
needed to regenerate the target, evidence that the representation is
information-complete rather than merely discriminative for retrieval. This
directly operationalizes the notion of semantic sufficiency that motivates DME.

\begin{table}[t]
\centering
\resizebox{\linewidth}{!}{%
\begin{tabular}{l ccc ccc ccc ccc}
\toprule
\multirow{2}{*}{Direction} & \multicolumn{3}{c}{Image} & \multicolumn{3}{c}{Video} & \multicolumn{3}{c}{VisDoc} & \multicolumn{3}{c}{All} \\
\cmidrule(lr){2-4}\cmidrule(lr){5-7}\cmidrule(lr){8-10}\cmidrule(lr){11-13}
 & @1 & @5 & @10 & @1 & @5 & @10 & @1 & @5 & @10 & @1 & @5 & @10 \\
\midrule
\multicolumn{13}{l}{\emph{Self (reconstruction completeness)}} \\
q2q & 0.8758 & 0.9469 & 0.9542 & 0.8781 & 0.8984 & 0.9072 & 0.8909 & 0.9041 & 0.9102 & 0.8792 & 0.9184 & 0.9262 \\
d2d & 0.7737 & 0.9117 & 0.9387 & 0.6431 & 0.8640 & 0.9136 & 0.9455 & 0.9805 & 0.9903 & 0.7432 & 0.9014 & 0.9355 \\
\midrule
\multicolumn{13}{l}{\emph{Cross (counterpart completeness)}} \\
q2d & 0.6279 & 0.8102 & 0.8568 & 0.5920 & 0.8328 & 0.8787 & 0.9513 & 0.9925 & 0.9987 & 0.6585 & 0.8453 & 0.8859 \\
d2q & 0.8736 & 0.9478 & 0.9551 & 0.8728 & 0.8989 & 0.9104 & 0.8978 & 0.9200 & 0.9337 & 0.8771 & 0.9216 & 0.9318 \\
\bottomrule
\end{tabular}}
\caption{\textbf{Representation completeness measured by teacher-forced Top-$K$
accuracy (acc@$K$).} Given a pooled, un-normalized embedding as the only
conditioning prefix, the trained DME model is asked to recover the target
tokens under teacher forcing; acc@$K$ is the fraction of evaluated positions
whose ground-truth token falls within the model's Top-$K$ predictions. Scores
are micro-averaged over the first $L_i$ evaluated positions of each sample ($\sum\text{hits}/\sum\text{positions}$),
and ``All'' merges the three domains with token-level weighting. Self directions
(q2q, d2d) probe how faithfully an embedding re-encodes its own text; cross
directions (q2d, d2q) probe how much of the paired counterpart it retains.}
\label{tab:ntp-topk-micro}
\end{table}

Table~\ref{tab:ntp-topk-micro} reports representation completeness across the
four directions and three domains. Overall, DME embeddings are highly
information-complete: on the merged set, the self directions recover the
ground-truth token at Top-1 in $87.9\%$ (q2q) and $74.3\%$ (d2d) of positions,
rising above $92\%$ at Top-10. This confirms that a single pooled vector
retains most of the token-level content of its own input rather than collapsing
to a coarse, retrieval-only summary. The cross directions are naturally harder,
since reproducing the counterpart requires information that is only shared
through the matched pair; nonetheless, d2q reaches $87.7\%$ at Top-1 and q2d
reaches $65.9\%$, indicating that a substantial fraction of the paired content
is genuinely encoded, consistent with the counterpart-side supervision used
during Stage-2 training. Across domains, VisDoc is the easiest to reconstruct
(q2d reaches $95.1\%$ at Top-1), matching its text-heavy, highly regular
content, whereas video is the hardest in the cross setting (q2d $59.2\%$ at
Top-1), reflecting the larger information gap between a text query and a
temporally extended visual target. Taken together, these results provide a
direct, quantitative confirmation that DME representations are semantically
sufficient, and the metric offers an interpretable signal that we further use
to guide representation optimization in industrial settings.


\subsection{Efficiency and Query Latency}
\label{sec:efficiency-query-latency}

A key design goal of DME is to introduce evidence-grounded latent reasoning without sacrificing the serving efficiency required by large-scale retrieval systems such as Douyin. Unlike explicit CoT generation or reranking-based reasoning, Stage 2-A only inserts a small number of anchor and typed latent tokens into the query encoder forward pass. We therefore measure the additional query-side latency introduced by these latent tokens.

The measurement only covers the forward pass of the query encoder. It excludes candidate encoding, query--candidate scoring, ANN retrieval, data loading, collation, and host-to-device transfer. 
We compare the same batches under two settings: \textbf{w/o Latent Tokens}, where the query uses the standard readout, and \textbf{w/ Latent Tokens}, where anchor and typed latent reasoning tokens are enabled. We use 8 high-performance GPUs with batch size 4. We first run 20 warmup iterations, which are not included in the timing statistics, and then record up to 200 measured iterations. Each iteration corresponds to one query-batch forward pass. We report p50 latency, \ie, the median latency over measured batches, which is a standard percentile metric for serving-time latency analysis.

\begin{table}[t]
\centering
\small
\setlength{\tabcolsep}{5pt}
\renewcommand{\arraystretch}{1.08}
\begin{tabular}{llcccc}
\toprule
\textbf{Query Type} & \textbf{Dataset}
& \shortstack{\textbf{w/o Latent}\\\textbf{p50 (ms/batch)}}
& \shortstack{\textbf{w/ Latent}\\\textbf{p50 (ms/batch)}}
& \shortstack{\textbf{$\Delta$}\\\textbf{(ms/batch)}}
& \shortstack{\textbf{$\Delta$}\\\textbf{(ms/query)}} \\
\midrule
Text & MSCOCO-T2I
& 41.962 & 45.450 & +3.487 & +0.87 \\
Image+Text & OK-VQA
& 111.084 & 111.399 & +0.315 & +0.08 \\
Video+Text & Video-MME
& 1049.479 & 1052.694 & +3.215 & +0.80 \\
\bottomrule
\end{tabular}
\caption{
\textbf{Query encoder p50 latency with and without latent reasoning tokens.} P50 denotes the median latency over measured batches. Per-query overhead is computed using batch size 4.
}
\label{tab:query-latency}
\end{table}

As shown in Table~\ref{tab:query-latency}, enabling latent reasoning introduces only minor overhead. The additional cost is below 1 ms per query for text and video queries, and nearly negligible for image--text queries. This is consistent with the design of DME: latent reasoning does not perform autoregressive generation or multi-round inference, but only adds a few soft tokens to the same encoder forward pass. For multimodal inputs, especially video, the computation is dominated by the original visual tokens, making the relative cost of the additional latent tokens small.

These results show that DME can keep latent reasoning active during query encoding while remaining close to the standard dense retrieval pipeline, which is important for practical deployment in large-scale multimodal retrieval.

